%% chapters/intro.tex
%%

% The \phantomsection command is needed to create a link to a place in the document that is not a
% figure, equation, table, section, subsection, chapter, etc.
% https://tex.stackexchange.com/questions/44088/when-do-i-need-to-invoke-phantomsection
\phantomsection

\chapter{\lang{Introduction}{Introdução}}
\phantomsection

As redes Peer-to-peer estão na computação como conceito e aplicações desde os
primórdios da ARPANET. Na verdade, a própria ARPANET poderia ser descrita como uma
rede de comunicação entre nodos - ou “pares” - descentralizada (Crocker, 1969), o
que se refletiu na própria criação da internet. A internet como conhecemos hoje é
uma comunicação entre múltiplas máquinas, múltiplos pares, para envio e recebimento
de recursos. Longe de ser uma novidade recente, o conceito é maduro e fundamenta a
comunicação entre sistemas desde a década de 1960.

O conceito evoluiu ao longo de três grandes gerações de redes peer-to-peer, sempre
adicionando ou melhorando a solução anterior. Na primeira geração houve a primeira
grande implementação dessa arquitetura para resolver o problema de transferência de
mídias digitais, ainda com servidores centrais direcionando o tráfego com tabelas de
índices. A segunda geração verdadeiramente descentralizou a rede, não dependendo mais
de nenhum servidor central. Então, a terceira geração de redes eliminou os principais
problemas de flooding, introduziu o upload em chunks e muitos outros problemas de suas
antecessoras, se tornando um sucesso global.

Entretanto, essa evolução se restringiu ao envio de arquivos estáticos. Arquivos nunca
mudam entre um envio e outro, uma música sempre será o mesmo file, terá o mesmo hash, os
mesmos metadados. Para as aplicações do mundo moderno, pautadas no envio e
recebimento massivo de dados de aplicações em tempo real, as redes peer-to-peer
enfrentam um novo grande desafio: a sincronização de dados em tempo real.

Como uma alternativa para superar essas limitações sem recorrer ao modelo
cliente-servidor, este trabalho propõe o desenvolvimento e análise de um
Motor de Sincronização Descentralizada para Aplicações em Tempo Real. Para
fundamentar e resolver o conflito de divergência sem autoridade central, o
motor utiliza estruturas de dados específicas para consistência eventual,
como Conflict-free Replicated Data Types (CRDTs) \cite{shapiro2011}. Para
validar e avaliar o desempenho da solução em termos de latência,
resiliência e concorrência, é desenvolvida uma Proof of Concept (PoC)
focada na sincronização espacial de mapas interativos entre múltiplos
usuários simultâneos.

\section{Objetivos}

Esta seção busca formalizar os objetivos do presente trabalho.

\subsection{Objetivos gerais}

Este trabalho tem como objetivo projetar, implementar e avaliar
experimentalmente um motor de sincronização de estado descentralizado
(Peer-to-Peer), de baixo nível e agnóstico a domínio, visando garantir a
Consistência Eventual Forte (SEC) em aplicações de tempo real por meio da
integração de Tipos de Dados Replicados Livres de Conflitos (CRDTs) e
ordenação lógica de eventos.

\subsection{Objetivos específicos}

\begin{itemize}
    \item Estudar os fundamentos teóricos de consistência em sistemas
    distribuídos descentralizados, analisando os conceitos de Consistência
    Eventual Forte (SEC), Tipos de Dados Replicados Livres de Conflitos
    (CRDTs), relógios lógicos e protocolos de fofoca (gossip);

    \item Analisar os requisitos técnicos e desafios práticos para o
    desenvolvimento de sistemas de alta performance em tempo real, mapeando
    os requisitos de paralelismo, controle de concorrência, latência e
    segurança no ciclo de vida da memória;

    \item Projetar a arquitetura modular de baixo nível do motor de
    sincronização, definindo os fluxos de entrada e saída de rede
    assíncrona, enquadramento (framing) de dados, controle de concorrência
    multi-thread e o modelo de persistência volátil;

    \item Implementar o middleware (motor de sincronização) focando na
    performance de execução, empregando mecanismos eficientes de
    serialização de dados (zero-copy) e estruturas de dados otimizadas
    para manipulação concorrente;

    \item Desenvolver um protótipo funcional de tempo real (aplicação de
    prova de conceito) que consuma as interfaces expostas pelo motor,
    demonstrando na prática o comportamento dinâmico e a natureza
    agnóstica do sistema;

    \item Avaliar a solução proposta por meio de testes experimentais,
    analisando métricas de tempo de convergência de réplicas, taxa de
    transferência de rede (payload de fofoca), latência induzida e
    consumo de recursos de processamento e memória;

    \item Documentar as limitações de design identificadas, os
    aprendizados de engenharia de software aplicados a sistemas
    distribuídos e as futuras direções para evolução da infraestrutura
    proposta.
\end{itemize}

\section{Escopo do trabalho}

\section{Método de pesquisa}

Este trabalho utiliza uma pesquisa aplicada, com foco em investigar um
problema prático da engenharia de software distribuída: a viabilidade e os
desafios de manter estados consistentes em redes de tempo real
descentralizadas sem coordenação síncrona global. A abordagem adotada é
experimental, baseada no projeto, desenvolvimento e avaliação empírica de
um protótipo funcional.

Inicialmente, será realizado um estudo teórico aprofundado a respeito das
bases matemáticas de CRDTs (com foco em semilattices superiores e mutações
baseadas em delta), relógios de Lamport, relógios vetoriais e algoritmos de
disseminação epidêmica (gossip protocols). Além disso, serão revisadas as
bibliotecas e técnicas modernas de programação de sockets de rede
assíncronos e de alta concorrência na linguagem de programação que será
adotada.

Em seguida, será modelada e documentada a arquitetura lógica do motor de
sincronização. Esse projeto envolverá o mapeamento das dependências entre
os módulos de rede (camada de transporte), de persistência do estado da
réplica, de ordenação lógica de eventos e de difusão de atualizações.

Após a etapa de projeto arquitetural, será desenvolvido o motor de
sincronização (middleware) sob uma estrutura multi-thread otimizada,
implementando as classes e interfaces correspondentes aos CRDTs e relógios
lógicos de forma genérica. Concomitantemente, será desenvolvida uma
aplicação de prova de conceito (como um simulador distribuído simples ou
um quadro colaborativo interativo) cuja única responsabilidade será
manipular o estado através do middleware, validando sua natureza agnóstica
a domínio.

Por fim, serão conduzidos testes práticos e ensaios experimentais sob
diferentes cenários de carga de trabalho e concorrência induzida. Esses
ensaios buscarão caracterizar quantitativamente a performance da solução
proposta, avaliando as métricas de tempo total de convergência do estado
de todas as réplicas ativas, taxa de pacotes descartados por falta de
causalidade, consumo de banda de rede (bits/segundo) e consumo de
processador e memória RAM nos nós da rede simulada. Os resultados obtidos
guiarão a redação das conclusões e limitações técnicas do sistema.

\section{O uso de Inteligência Artificial neste trabalho}

Em conformidade com as boas práticas de transparência acadêmica,
esclarece-se nesta seção como ferramentas de Inteligência Artificial (IA)
foram empregadas na elaboração deste trabalho.

Um assistente de IA (Claude, da Anthropic, operando através da ferramenta
Claude Code) foi utilizado como apoio em tarefas de natureza instrumental e
editorial, e não na concepção científica da pesquisa. Especificamente, a
ferramenta auxiliou:

\begin{itemize}
    \item Na migração do texto deste documento, originalmente redigido pelo
    autor em processador de texto, para o template \LaTeX{} da UFSC em
    conformidade com as normas ABNT, preservando o conteúdo e a redação
    originais do autor;

    \item Na configuração do ambiente de desenvolvimento \LaTeX{} local
    (distribuição TeX Live, editor, extensões) e do fluxo de versionamento
    do projeto via \texttt{git};

    \item Na formatação das referências bibliográficas citadas neste
    capítulo no formato BibTeX e na inserção dos respectivos comandos de
    citação no texto;

    \item Na organização geral do repositório do projeto, incluindo a
    remoção de conteúdo de demonstração herdado do template original.
\end{itemize}

A definição do tema de pesquisa, a revisão da literatura, a redação do
conteúdo técnico e científico, as decisões metodológicas e as conclusões
apresentadas ao longo deste trabalho são de autoria e responsabilidade do
autor. O uso da IA restringiu-se, portanto, a um papel de ferramenta de
produtividade na produção do documento, análogo ao uso de um editor de
texto ou de um verificador ortográfico, sem substituir o julgamento
científico ou a autoria intelectual da pesquisa.

\section{Organização do trabalho}

Além desta introdução, o restante deste trabalho está organizado da
seguinte forma: o capítulo de Desenvolvimento apresenta a fundamentação
teórica sobre consistência em sistemas distribuídos descentralizados --
CRDTs, relógios lógicos e protocolos de disseminação epidêmica -- seguida
do projeto da arquitetura modular do motor de sincronização e dos detalhes
de sua implementação, com foco em desempenho e concorrência. Em seguida, é
apresentado o protótipo de prova de conceito desenvolvido sobre o motor, e
a avaliação experimental conduzida sobre ele, caracterizando métricas de
convergência, latência, tráfego de rede e consumo de recursos sob
diferentes cenários de carga. Por fim, o capítulo de Conclusão retoma os
objetivos definidos neste capítulo, discute os resultados obtidos e as
limitações identificadas, e aponta direções para trabalhos futuros.
